% problem-000218 11 羰基化合物多步反应
\section*{11. 羰基化合物多步反应}
\textit{下列为分问。}

\subsection*{11-1}
α-杂原⼦取代羰基化合物在天然产物和药物分⼦的合成中具有重要⽤途。酰胺化合物A在三氟甲磺酸酐(Tf_{2}O)和2-碘吡啶(2-I-py)存在下，与⼿性亚磺酰胺试剂反应，⽣成α-氨基酰胺化合物H。反应过程如下：

（图：）

\subsection*{11-1}
-1 画出中间体B、C、F和G的结构。

\subsection*{11-1}
-2 中间体D和E哪个形成更快？说明原因。

\subsection*{11-1}
-3 画出G到H的过渡态结构。

\subsection*{11-2}
如下过程可⽤于α-羟基酮的制备。该反应的关键在于中间体L的形成以及L中羰基在光激发下形成双⾃由基M。

（图：）

\subsection*{11-2}
-1 画出 L 的结构。

\subsection*{11-2}
-2 中间体M和N都是双⾃由基，画出M、N和O的结构。

\subsection*{11-2}
-3 采⽤亚磷酸三异丙酯R作为反应底物时，可⽣成与O类似的O1中间体。画出⽤结构简式表述的O1到S的反应机理。

（图：）

\subsection*{11-2}
-4 画出化合物S经硼氢化钠处理后得到的最终产物结构。
